\section{Outlook}

A direct extension of the inclined model would be a continuous connection between the shoulder and peak. The existing inclined samples provide a basis for testing such a modification without changing the detector simulation. Independently of this improvement, the current models could serve as a starting point for beam reconstruction. This would first require establishing how their shape parameters depend jointly on gamma energy and incidence angle. Tests under beam conditions not used to establish this relationship could then determine which beam parameters can be recovered from spatial shape alone and where different conditions remain indistinguishable.

Further simulations could test how the observed behaviour extends beyond the present crystal geometry and energy range. Varying the crystal thickness would change both the available interaction depths and the optical propagation distances, allowing their influence on the central peak and broad component to be studied systematically. Extending the gamma energy range above the pair-production threshold would also test how the additional charged-particle tracks and annihilation photons affect the optical distribution.

Experimental validation would be needed to establish whether the simulated dependence of the optical distribution on gamma energy and incidence angle is also observed in a physical detector. A comparison with measurements could test the predicted central peak and tails and help assess the optical-property assumptions. This would require accounting for the response of the detector used in the comparison.